\section{PERFIL PROFISSIONAL DO EGRESSO}

Nos termos da \textbf{Resolução CNE/CES nº 2, de 24 de abril de 2019}, e da \textbf{Resolução CNE/CES nº 1, de 26 de março de 2021}, que instituem as Diretrizes Curriculares Nacionais para os cursos de Graduação em Engenharia, observa-se, especialmente, o disposto no Art. 3º, o qual estabelece que o perfil do egresso deve contemplar determinadas competências e características fundamentais à formação do engenheiro, a saber:

\begin{adjustwidth}{4cm}{0cm}
\it \footnotesize \setstretch{.5}

	O perfil do egresso do curso de graduação em Engenharia deve compreender, entre outras, as seguintes características:
	
	\begin{enumerate}[label=\Roman* -]
		\item visão holística e humanista, postura crítica, reflexiva, criativa, cooperativa e ética, sustentada por sólida formação técnico-científica;
		\item capacidade para pesquisar, desenvolver, adaptar e empregar novas tecnologias, atuando de modo inovador e empreendedor;
		\item habilidade para reconhecer necessidades dos usuários, formular, analisar e solucionar problemas de Engenharia de forma criativa;
		\item adoção de perspectivas multidisciplinares e transdisciplinares na prática profissional;
		\item consideração de aspectos globais, políticos, econômicos, sociais, ambientais, culturais, de segurança e de saúde no trabalho;
		\item atuação isenta e comprometida com a responsabilidade social e o desenvolvimento sustentável.
	\end{enumerate}
	
\end{adjustwidth}

Nesse sentido, a partir dessas diretrizes, a pergunta-guia que estrutura a definição das competências do egresso, nos termos do \cite{res259}, é a seguinte: \textit{Quais conhecimentos, habilidades e atitudes o estudante de engenharia deve possuir ao concluir o curso e em que nível de proficiência?}

Diante dessa indagação e em consonância direta com os princípios estabelecidos pelas Diretrizes Curriculares Nacionais, estabelece-se que, ao finalizar sua formação, o futuro \NOMEPROFISSIONAL\ deverá evidenciar conhecimentos, habilidades e atitudes compatíveis com o exercício responsável e tecnicamente qualificado da Engenharia, alcançando níveis de proficiência que assegurem autonomia profissional, rigor científico e comprometimento ético. Tal perfil compreende:

\begin{itemize}
	\item \textbf{Conhecimentos}: domínio consistente dos fundamentos matemáticos, científicos e tecnológicos da Engenharia, aplicados com precisão conceitual e capacidade analítica suficiente para sustentar modelagens, diagnósticos e decisões técnicas de maneira autônoma.
	
	\item \textbf{Habilidades}: aptidão para identificar demandas, analisar contextos e propor soluções de Engenharia com criatividade, rigor metodológico e uso competente de tecnologias contemporâneas; capacidade de pesquisar, desenvolver e adaptar soluções técnicas, atuando de forma integrada em equipes multidisciplinares, em nível de proficiência equivalente ao desempenho profissional inicial.
	
	\item \textbf{Atitudes}: postura ética, crítica e responsável, orientada pela consciência dos impactos sociais, econômicos, ambientais, culturais, políticos, de segurança e de saúde inerentes à prática da Engenharia; compromisso permanente com a responsabilidade social, a cooperação, a inovação e o desenvolvimento sustentável, demonstrando maturidade no processo de tomada de decisão.
\end{itemize}
